% ==============================================================================
% 1.9 a 1.12: ORGANIZACIÓN, SEGUIMIENTO, TECNOLOGÍAS Y ANEXOS
% ==============================================================================
\subsection{1.9. Organización interna}

El proyecto lo desarrolla NEXWARE, un equipo de desarrollo de software de dos integrantes, Mauro Boris Gutierrez Primintela y Carlos Leonel Barrios Mayorga, bajo la guía del docente de la asignatura, Ing. José Junior Villagómez Melgar. Con solo dos personas no se justifica una jerarquía de mando: ambos participan en las decisiones de diseño, comparten la responsabilidad por el resultado final y se revisan mutuamente el trabajo. Esta forma de trabajo corresponde a un equipo de tipo democrático descentralizado en la clasificación de equipos de software de \textcite{pressman2010ingenieria}.

El trabajo técnico se reparte en dos bloques de igual esfuerzo (58 puntos de dificultad cada uno), definidos de manera que cada integrante sea dueño de punta a punta (base de datos, interfaz de servicio, pantallas y pruebas) de los casos de uso que le corresponden. La Tabla~\ref{tab:organizacion_roles} resume el reparto.

\begin{table}[H]
\caption{Organización del equipo y reparto del trabajo}
\label{tab:organizacion_roles}
\begin{threeparttable}
\begin{tabular}{>{\raggedright\arraybackslash}p{2.6cm} >{\raggedright\arraybackslash}p{5cm} >{\raggedright\arraybackslash}p{5.2cm}}
\toprule
Aspecto & Mauro Boris Gutierrez Primintela (bloque A) & Carlos Leonel Barrios Mayorga (bloque B) \\
\midrule
Bloques de casos de uso & Cuentas y acceso, auditoría y adjudicación, reportes de Business Intelligence, notificaciones y consulta del estado de la postulación & Convocatoria y postulación, preselección con inteligencia artificial, becados y asistencia, y consulta de bitácora \\
Casos de uso & 13 & 13 \\
Tareas adicionales (no son casos de uso) & 3 & 1 \\
Puntos de dificultad & 58 & 58 \\
Plataformas principales & Web (React) e inicio de sesión y perfil móvil & Web, móvil (Flutter) y módulo de inteligencia artificial (Python) \\
\bottomrule
\end{tabular}
\begin{tablenotes}
\item \textit{Nota.} Elaboración propia, a partir del plan de desarrollo del proyecto. La integración entre ambos bloques se limita a contratos de datos acordados al inicio.
\end{tablenotes}
\end{threeparttable}
\end{table}

Las decisiones que afectan a ambos bloques (esquema de la base de datos, estados de la postulación y contratos de integración) se toman de común acuerdo antes de que cada integrante se concentre en su bloque. NEXWARE es la empresa desarrolladora y la DUBSS es el cliente. La estructura es plana: Mauro Boris Gutierrez Primintela y Carlos Leonel Barrios Mayorga son líderes en igualdad de condiciones, cada uno responsable de su bloque, y ambos rinden cuentas al docente guía. La Figura~\ref{fig:organigrama} presenta el organigrama.

\begin{figure}[H]
\caption{Organigrama de NEXWARE}
\label{fig:organigrama}
\centering
\begin{tikzpicture}[node distance=1cm, box/.style={draw, rounded corners, align=center, minimum width=5cm, minimum height=1.1cm}]

\node[box] (cli) {Cliente: DUBSS};

\node[box, below=of cli] (neo) {NEXWARE\\Dirección compartida};

\node[box, below left=of neo, xshift=-0.3cm] (a) {Líder de bloque A\\Mauro Boris Gutierrez Primintela};

\node[box, below right=of neo, xshift=0.3cm] (b) {Líder de bloque B\\Carlos Leonel Barrios Mayorga};
\draw (cli) -- (neo); \draw (neo) -- (a); \draw (neo) -- (b);
\end{tikzpicture}

\smallskip

\noindent\raggedright\textit{Nota.} Elaboración propia.
\end{figure}

\subsection{1.10. Mecanismos de seguimiento y control}

El seguimiento del proyecto se apoya en el plan de desarrollo, que descompone el sistema en 31 tareas (27 tareas implementan los 26 casos de uso, ya que la precalificación con inteligencia artificial se divide en dos, y cuatro son tareas técnicas adicionales: tres del bloque A y una del B) y en seis hitos de integración (H1 a H6), el último de los cuales es la prueba de extremo a extremo local del flujo completo: postular, puntuar, priorizar, auditar, adjudicar, registrar asistencia y consultar indicadores. Cada tarea tiene un responsable, una estimación en puntos y una definición de terminado (base de datos, servicio, pantallas y pruebas).

El control se realiza mediante tres formularios de reporte que cada integrante completa por tarea y que se revisan en las reuniones de seguimiento del equipo. Las Tablas~\ref{tab:form_finalizadas}, \ref{tab:form_proceso} y \ref{tab:form_errores} describen sus campos. Los formularios se registran junto al repositorio del código: cada tarea tiene una ficha de unidad de trabajo (carpeta \texttt{docs/UnidadesTrabajo}) con su alcance, su dependencia y los pasos de validación; el estado de cada caso de uso (Pendiente, En progreso o Hecho) se mantiene en su propio archivo de caso de uso, y los cambios se integran mediante \textit{pull requests} de GitHub que revisa el equipo. El equipo se reúne una vez por semana para revisar el avance, los impedimentos y los errores abiertos. Además, se realiza una breve coordinación diaria por mensajería cuando existen dependencias entre los dos bloques.

Además de los formularios, el repositorio aplica controles automáticos en cada cambio. Un \textit{hook} de Git valida que cada \textit{commit} siga el formato \texttt{tipo[ALCANCE]: descripción} (herramienta commitlint) y que el código modificado cumpla el estilo de su lenguaje (ESLint, Checkstyle y \textit{flutter analyze}). Luego, GitHub Actions ejecuta, solo para los proyectos que cambiaron, el análisis de estilo, las pruebas automatizadas y la compilación de cada uno de los cinco proyectos (backend principal, módulo de inteligencia artificial, panel web, landing y aplicación móvil). En el backend, las pruebas de integración usan una base PostgreSQL real, de modo que las migraciones de la base de datos también se verifican en cada cambio.

\begin{table}[H]
\caption{Formulario de reporte de tareas finalizadas}
\label{tab:form_finalizadas}
\begin{threeparttable}
\begin{tabular}{>{\raggedright\arraybackslash}p{4cm} >{\raggedright\arraybackslash}p{9cm}}
\toprule
Campo & Contenido \\
\midrule
Tarea y caso de uso & Identificador de la tarea (por ejemplo, A9) y caso de uso que implementa. \\
Responsable y fecha & Integrante que la completó y fecha de cierre. \\
Puntos estimados y esfuerzo real & Puntos asignados en el plan y días-persona efectivamente invertidos. \\
Entregables & Base de datos, servicio, pantallas y pruebas terminados. \\
Criterios de aceptación & Resultado de la verificación contra los criterios del requisito funcional asociado. \\
Contratos afectados & Contratos de integración que la tarea produce o consume. \\
\bottomrule
\end{tabular}
\begin{tablenotes}
\item \textit{Nota.} Elaboración propia.
\end{tablenotes}
\end{threeparttable}
\end{table}

\begin{table}[H]
\caption{Formulario de reporte de tareas en proceso}
\label{tab:form_proceso}
\begin{threeparttable}
\begin{tabular}{>{\raggedright\arraybackslash}p{4cm} >{\raggedright\arraybackslash}p{9cm}}
\toprule
Campo & Contenido \\
\midrule
Tarea y responsable & Identificador de la tarea y del integrante que la desarrolla. \\
Porcentaje de avance & Avance declarado respecto de la definición de terminado. \\
Trabajo restante & Actividades pendientes y fecha estimada de cierre. \\
Dependencias & Tareas o contratos del otro integrante de los que depende. \\
Impedimentos & Obstáculos que requieren una decisión o ayuda del equipo. \\
\bottomrule
\end{tabular}
\begin{tablenotes}
\item \textit{Nota.} Elaboración propia.
\end{tablenotes}
\end{threeparttable}
\end{table}

\begin{table}[H]
\caption{Formulario de reporte de errores}
\label{tab:form_errores}
\begin{threeparttable}
\begin{tabular}{>{\raggedright\arraybackslash}p{4cm} >{\raggedright\arraybackslash}p{9cm}}
\toprule
Campo & Contenido \\
\midrule
Identificador y fecha & Número correlativo del error y fecha de detección. \\
Caso de uso y componente & Caso de uso afectado y componente (web, móvil, servicio o módulo de inteligencia artificial). \\
Descripción y pasos & Descripción del fallo y pasos para reproducirlo. \\
Severidad & Crítico, alto, medio o bajo, según el impacto sobre el flujo de la postulación. \\
Responsable y estado & Integrante asignado; estado abierto, en corrección o cerrado. \\
\bottomrule
\end{tabular}
\begin{tablenotes}
\item \textit{Nota.} Elaboración propia.
\end{tablenotes}
\end{threeparttable}
\end{table}

\subsection{1.11. Tecnologías para el desarrollo de proyecto}

Las tecnologías del sistema son las que figuran en los archivos de dependencias del repositorio; la Tabla~\ref{tab:tecnologias_reales} las resume por componente. El sistema se organiza como un repositorio único (\textit{monorepo}) con cinco proyectos independientes: el backend principal, el módulo de inteligencia artificial, el panel web, la landing institucional y la aplicación móvil, más la configuración de infraestructura con la que se levanta todo mediante Docker Compose. La justificación de estas tecnologías y su descripción detallada se presentan en el Capítulo 6 (Infraestructura para la Producción de Software); las herramientas y el equipo físico del grupo figuran en la sección 1.7 de este capítulo.

\begin{longtable}{>{\raggedright\arraybackslash}p{2.9cm} >{\raggedright\arraybackslash}p{5.6cm} >{\raggedright\arraybackslash}p{4.5cm}}
\caption{Tecnologías utilizadas en el desarrollo del sistema}\label{tab:tecnologias_reales}\\
\toprule
Componente & Tecnologías y versiones & Uso en el sistema \\
\midrule
\endfirsthead
\toprule
Componente & Tecnologías y versiones & Uso en el sistema \\
\midrule
\endhead
\bottomrule
\endfoot
Backend principal & Java 21; Spring Boot 4.1.2 (Web MVC, Data JPA, Security con servidor de recursos JWT, Validation, Mail, Actuator y AspectJ); Maven & API REST (\texttt{/api/v1}) con arquitectura por módulos y capas, autenticación propia con tokens JWT, auditoría en bitácora y envío de correos \\
Persistencia & PostgreSQL 18; Flyway; Hibernate (JPA); Redis; MinIO & Base de datos con migraciones versionadas, caché y almacenamiento de documentos y actas en PDF \\
Servicios del backend & Bucket4j; springdoc-openapi (Swagger UI); Firebase Admin (FCM); Micrometer con Prometheus; Logback con salida JSON; Lombok & Limitación de peticiones por IP, documentación de la API, notificaciones push, métricas y registros estructurados \\
Panel web & React 19; TypeScript 6; Vite 8; TanStack Query, Router, Form y Table; Tailwind CSS 4 con shadcn; Zod 4; Zustand 5; Axios; Recharts & Interfaz de los roles administrativos y de los postulantes, con tablas paginadas, formularios validados y gráficos del tablero de indicadores \\
Aplicación móvil & Flutter 3.44 con Dart 3.12; Riverpod 2.6; go\_router 18; dio 5; flutter\_secure\_storage; Firebase Messaging & Aplicación Android del postulante: acceso, perfil, consulta de la postulación y notificaciones push \\
Módulo de inteligencia artificial & Python 3.12; FastAPI; Uvicorn; pytest & Servicio independiente para el cálculo del puntaje de vulnerabilidad \\
Landing institucional & Astro 7 con TypeScript; Vitest & Sitio estático bilingüe (español e inglés) de la empresa NEXWARE \\
Infraestructura de desarrollo & Docker Compose; nginx; Mailpit; Prometheus; Grafana & Entorno completo reproducible en cada equipo y observabilidad del backend \\
Calidad e integración continua & GitHub Actions; Husky, lint-staged y commitlint; ESLint y Prettier; Checkstyle; Ruff; JUnit con Testcontainers; Vitest; flutter test & Verificación automática del estilo, de las pruebas y de la compilación de cada proyecto \\
\end{longtable}
\noindent\textit{Nota.} Elaboración propia a partir de los archivos de dependencias y de configuración del repositorio.


\subsection{1.12. Anexos}

\subsubsection{1.12.1. Líneas de código de Software de referencia 1: TastyIgniter}

Salida de \textit{cloc} 2.12 sobre el repositorio clonado (primero el repositorio TastyIgniter y luego el paquete tastyigniter/core).

\begingroup\footnotesize
\begin{verbatim}
github.com/AlDanial/cloc v 2.12  T=1.23 s (53.6 files/s, 2622.3 lines/s)
-------------------------------------------------------------------------------
Language                     files          blank        comment           code
-------------------------------------------------------------------------------
PHP                             54            405           1198           1198
JSON                             2              0              0            120
Markdown                         2             25              0             62
YAML                             2              9              0             62
Bourne Shell                     1              8              0             55
XML                              1              0              2             29
SVG                              1              0              0             18
INI                              1              4              0             14
JavaScript                       1              2             11              4
Text                             1              0              0              2
-------------------------------------------------------------------------------
SUM:                            66            453           1211           1564
-------------------------------------------------------------------------------

github.com/AlDanial/cloc v 2.12  T=26.80 s (54.7 files/s, 6251.3 lines/s)
-------------------------------------------------------------------------------
Language                     files          blank        comment           code
-------------------------------------------------------------------------------
PHP                           1030          21860          13300          84276
JSON                            17              0              0          20671
JavaScript                     138           1472           1292          14192
Blade                          218            212              1           5614
SCSS                            36            342            131           2688
CSS                             20            181            122            926
YAML                             1              8              0             70
Markdown                         2             16              0             42
Text                             1              5              0             37
XML                              1              0              0             28
SVG                              1              0              0             18
INI                              1              4              0             14
-------------------------------------------------------------------------------
SUM:                          1466          24100          14846         128576
-------------------------------------------------------------------------------
\end{verbatim}
\endgroup

\subsubsection{1.12.2. Líneas de código de Software de referencia 2: Floreant POS}

Salida de \textit{cloc} 2.12 sobre el repositorio clonado.

\begingroup\footnotesize
\begin{verbatim}
github.com/AlDanial/cloc v 2.12  T=16.92 s (68.1 files/s, 14309.1 lines/s)
----------------------------------------------------------------------------------------
Language                              files          blank        comment           code
----------------------------------------------------------------------------------------
Java                                    938          28031          33334         102268
Text                                      2              1              0          42622
Jasper Report XML/Template               34             12              4          14689
XML                                      72            616             21          11548
Properties                                9            363              4           4872
Hibernate                                81            241             76           1843
Maven                                     5             40              6            656
HTML                                      1              8             11            593
Markdown                                  1             46              0            131
DOS Batch                                 9              1              0             56
Bourne Shell                              1              2              0              1
----------------------------------------------------------------------------------------
SUM:                                   1153          29361          33456         179279
----------------------------------------------------------------------------------------
\end{verbatim}
\endgroup

\subsubsection{1.12.3. Líneas de código de Software de referencia 3: OSPOS}

Salida de \textit{cloc} 2.12 sobre el repositorio clonado.

\begingroup\footnotesize
\begin{verbatim}
github.com/AlDanial/cloc v 2.12  T=25.01 s (56.5 files/s, 6185.7 lines/s)
-------------------------------------------------------------------------------
Language                     files          blank        comment           code
-------------------------------------------------------------------------------
PHP                           1289          10951          11357         115287
JSON                             3              0              0           6264
SQL                             40            605            468           1867
YAML                            19            199            106           1486
Markdown                        12            387              4           1393
CSS                             12            235             38           1221
JavaScript                       6            158            354            742
Bourne Shell                     6            109             53            465
SVG                              9              0              0            422
PowerShell                       3             74             33            170
HTML                             8             16              0             72
XML                              1              0              4             51
Dockerfile                       2             10              0             36
TOML                             1              1              2             25
INI                              1              3              0             12
CSV                              1              0              0              2
Text                             1              0              0              2
-------------------------------------------------------------------------------
SUM:                          1414          12748          12419         129517
-------------------------------------------------------------------------------
\end{verbatim}
\endgroup

